\documentclass[review]{elsarticle}
% Compile with: pdflatex si_middleware_extension.tex (run twice for cross-references)
% "review" option = double-spaced, numbered lines, single column --- standard for JSS submission.
%
% This file is the Elsevier-formatted counterpart of si_middleware_extension.md (the source of
% record for content). It extends the rejected Middleware 2026 submission per the revision plan
% in ../middleware2026/middleware26_revision_plan.md, and reports two new experiments (S5.5, S5.6)
% run for this revision. See si_middleware_extension.md's front matter for status notes.

\usepackage[utf8]{inputenc}
\usepackage[T1]{fontenc}
\usepackage{amsmath}
\usepackage{amssymb}
\usepackage{hyperref}
\usepackage{booktabs}
\usepackage{array}
\usepackage{longtable}

\journal{Journal of Systems and Software}

\begin{document}

\begin{frontmatter}

\title{Heterogeneous Graph Learning for Cascade Impact Prediction in Distributed Publish-Subscribe Middleware}

\author{Author(s) TBD}
\address{Institution(s) TBD}

\begin{abstract}
Hardening publish-subscribe middleware against cascading service disruptions requires propagating
predicted failures prior to deployment. Dependencies among publishers, subscribers, topics,
brokers, libraries, and deployment nodes were represented by graph-based methods, but static
analysis methods failed to capture the typed semantics of complex, sparse, and decoupled
topologies required to model their runtime behavior. We present a heterogeneous graph learning
(HGL) framework for impact prediction of application-level faults without measured run-time data.
HGL acts on the architecture model across Applications, Libraries, Topics, Brokers, and deployment
Nodes, and learns a relation-specific message function for each publish-subscribe relation, with
QoS attributes directly attributed to edges as specified in the definition. Crucially, while our
model maps the full multi-tier infrastructure to serve as a rich contextual embedding layer, our
current prediction targets and validation focus on application-level cascade impacts. For
baselines, instead of consuming a type-collapsed logical dependency projection of the same system,
the value of preserving node and relation types is isolated. We evaluate HGL across seven
publish-subscribe scenarios using a controlled factorial design that varies the graph-learning
architecture and QoS encoding, as well as baseline structural and homogeneous graph learning. To
avoid circular validation in the absence of telemetry, ground-truth impact is obtained from a
discrete-event simulation, and the baselines' topology features are fully decoupled from the GNN.
HGL achieves the best mean global ranking correlation ($\rho=0.411$ vs.\ GL's 0.186) and
identification F1 (0.688 vs.\ 0.596), outperforming homogeneous GNNs within a trained scenario's
distribution. Leave-One-Scenario-Out cross-validation, evaluating generalization to a genuinely
unseen topology, tells a different story: as trained here, neither HGL nor HGL-QoS transfers to
held-out scenarios, both falling to negative mean rank correlation and underperforming even the
homogeneous baselines out-of-distribution. We report this as found rather than adjusted to fit the
in-distribution narrative, since it directly contradicts what in-distribution metrics alone would
suggest and exposes a real generalization gap. The retention of typed heterogeneous semantics is
therefore a substantive but conditional advantage: decisive within a trained scenario, but not, on
this evidence, sufficient by itself for zero-shot transfer to an unseen pub-sub topology.
\end{abstract}

\begin{highlights}
\item Heterogeneous graph learning (HGL) beats homogeneous GNN baselines in-distribution on typed pub-sub graphs, but not out-of-distribution.
\item Formal $Q^*(v)$/$I^*(v)$ definitions resolve a simulator-attribution ambiguity via code audit.
\item A reversed-projection ablation flips $\rho$ from $+0.84$ to $-0.79$, confirming the dependency direction.
\item A traced, undisclosed ensemble step explains why this revision's rerun reverses the conference submission's Leave-One-Scenario-Out result -- reported transparently rather than suppressed.
\end{highlights}

\begin{keyword}
publish-subscribe middleware \sep heterogeneous graph neural networks \sep cascading failure \sep
pre-deployment verification \sep quality of service \sep architectural dependability
\end{keyword}

\end{frontmatter}

% ===========================================================================
\section{Introduction}

Publish-subscribe middleware --- standardized by the Data Distribution Service (DDS) and MQTT and
studied since the foundational surveys of Eugster et al.~\cite{Eugster2003} and Carzaniga et
al.~\cite{Carzaniga2001} --- decouples producers and consumers in time, space, and synchronization
\cite{OMG2015DDS,OASIS2019MQTT}. This decoupling is what lets such systems scale, and it is also
what makes them hard to reason about: a component's failure does not stop at its direct neighbors
but propagates along publish/subscribe routing paths that are invisible to source-level inspection.
Guaranteeing reliability therefore requires identifying, before deployment, which components would
cause the largest cascade if they failed --- a question classical dependability research on error
and attack tolerance in complex networks \cite{Albert2000} and cascading failures in interdependent
networks \cite{Buldyrev2010} answers for generic graphs, but not for the typed, QoS-annotated
topology of pub-sub middleware specifically. Structural centrality measures such as betweenness
\cite{Freeman1977,Brandes2001} offer a cheap, interpretable first approximation, but they treat
every edge as identical, so they cannot distinguish a publication edge from a shared-library
dependency or a broker-routing edge --- distinctions that determine whether a failure cascades
sequentially or blasts out to many components at once. Recent learning-based approaches close part
of this gap: homogeneous graph neural networks (GNNs) such as GCN \cite{Kipf2017} and GraphSAGE
\cite{Hamilton2017}, and attention variants such as GAT \cite{Velickovic2018}, have been used to
predict critical nodes directly from graph structure (FINDER~\cite{Fan2020}, DrBC~\cite{Munikoti2022},
PowerGraph~\cite{Munikoti2024}). But by construction these operate on a single node and edge type,
so applying them to pub-sub middleware requires first collapsing Applications, Topics, Brokers,
Libraries, and Nodes --- and the five distinct relations that connect them --- into one
undifferentiated graph. That collapse is exactly what discards the semantic distinctions structural
centrality already misses.

To resolve these limitations, we introduce Heterogeneous Graph Learning (HGL), a framework that
harnesses heterogeneous graph attention networks
\cite{Schlichtkrull2018,Wang2019,Hu2020,Fu2020} to identify critical components before deployment.
HGL models the publish-subscribe system as a native heterogeneous graph consisting of Applications,
Libraries, Topics, Brokers, and deployment Nodes, all connected by typed transport relations---such
as publication, subscription, library usage, broker connectivity, and host deployment. The
framework learns a distinct message function for each relation type, helping the model to
distinguish between dependencies that are structurally similar but semantically different, and thus
have varying real-world failure impacts.

\textbf{Running example.} Consider a simplified excerpt of an air-traffic-management (ATM)
deployment (26 applications, 8 libraries, 27 topics, 5 brokers, 8 nodes --- outside the seven
scenarios evaluated in \S~4.6, used here purely as an illustration): a \texttt{flight-plan-processor}
application publishes to the \texttt{flight.plans} topic, which a \texttt{clearance-router} and a
\texttt{trajectory-predictor} both subscribe to; all three link against a shared
\texttt{icao-protocol} library. Losing the publisher degrades the two subscribers through a cascade
that unfolds as messages go undelivered --- the kind of propagation a discrete-event simulator can
trace step by step. Losing the shared library, by contrast, is a simultaneous blast: all three
applications fail at once, a mode that ordinary connectivity metrics do not distinguish from an
ordinary dependency edge, because they collapse \texttt{publishes-to}, \texttt{subscribes-to}, and
\texttt{uses} into the same kind of edge. \S~3.1 defines the seven typed relations that keep these
apart; \S~3.2 defines the 7-dimensional QoS encoding that further distinguishes, e.g., a
\texttt{CRITICAL}-priority conflict-alert topic from a \texttt{BEST\_EFFORT} telemetry feed.

A key feature of HGL is its ability to perform message passing across all five node types, rather
than flattening the graph or collapsing node and edge types into a single view. High-fan-out Topic
hubs and Broker layers receive type-specific attention instead of being subsumed under a shared
transformation. To assess the value of this heterogeneity, we compare HGL against a homogeneous
baseline using a \emph{logical dependency projection} of the same system: pub-sub routing is
abstracted into application-level \texttt{DEPENDS\_ON} edges (for example, an edge $B \to A$ is
added if Application $B$ subscribes to a Topic that Application $A$ publishes, because the
subscriber \emph{depends on} the publisher; similarly, an edge $A \to L$ is added if $A$ uses
Library $L$, because the application depends on the library), with all node and edge types
collapsed into a single graph. This comparison isolates the necessity of preserving architectural
semantics. In the absence of real-world runtime telemetry to thoroughly evaluate the correctness of
the models in our work, we use a discrete-event simulation approach across seven scenarios (\S~3.4
defines the ground truth precisely). This results in a dense, dynamic cascade target and reduces
the validation circularity problem that arises when topology-based features are used both as inputs
and as labels.

Beyond the structural and learning-based baselines above, prior work on pub-sub dependability has
largely targeted the runtime and protocol layers --- fault tolerance, reliable event dissemination,
and recovery after a failure is observed \cite{Eugster2003,Carzaniga2001} --- rather than
pre-deployment estimation of which components would matter most if they failed; \S~2.1 surveys this
line in full. Our work instead asks that question before the system exists in production, from the
architecture model alone.

Our main contributions of this work are as follows:

\begin{enumerate}
\item We cast pre-deployment critical-component identification as a heterogeneous graph learning
  problem over the system's native, typed architecture graph.
\item We show that our heterogeneous approach leads to considerable improvements in detecting
  critical components within a trained scenario's distribution (an F1-score increase of $+0.092$
  over homogeneous GL) and exceeds traditional homogeneous graph-learning baselines whilst
  maintaining stronger ranking performance ($\Delta\rho=+0.225$).
\item We find that these in-distribution gains come predominantly from modeling typed nodes and
  relations, not from explicit QoS encoding. Out-of-distribution, however, neither typed message
  passing nor QoS encoding is sufficient by itself: Leave-One-Scenario-Out evaluation shows both
  HGL and HGL-QoS falling to negative mean rank correlation on unseen scenarios, underperforming
  even the homogeneous baselines -- a generalization gap that a purely in-distribution evaluation
  would not reveal, which we report as an honest negative result rather than omit.
\end{enumerate}

The remainder of this paper is organized as follows. Section~2 reviews related work. Section~3
describes the model, graph representations, and model variants. Section~4 defines the assessment
metrics. Section~5 presents the evaluation results, including ablation experiments and a
generalization analysis using a leave-one-out approach on previously unobserved scenarios.
Section~6 discusses threats to validity, and Section~7 concludes.

% ===========================================================================
\section{Related Work}

This study builds upon four areas of prior research: the dependability of publish-subscribe
middleware, structural analysis methods for detecting critical system components, learning-based
approaches for predicting critical nodes, and recent advances in heterogeneous graph neural
networks. While each of these areas offers important approaches and analytical methods for handling
complex distributed systems, none directly addresses the challenge of pre-deployment cascade
prediction for typed, graph-oriented models specifically designed for publish-subscribe middleware.
Our work seeks to fill this gap through integrating and broadening these approaches in an integrated
framework.

\subsection{Publish-Subscribe Middleware and Dependability}

Dependability of pub-sub middleware has been studied chiefly at the protocol, broker-overlay, and
runtime-recovery levels --- reliable delivery guarantees, replication, and post-failure recovery
mechanisms --- rather than at the level of pre-deployment, architecture-only estimation of which
components matter most; this section situates our work against that literature before turning, in
\S~2.2--2.4, to the structural and learned criticality methods it builds on. The publish-subscribe
paradigm has been widely adopted as a core communication abstraction in large-scale distributed
systems. Seminal studies have shown that pub-sub architectures achieve strong decoupling of
producers and consumers across time, space, and synchronization dimensions \cite{Eugster2003}.
Complementary lines of research in content-based networking highlight the flexibility of event
routing and the significance of efficient subscription matching in brokered overlays
\cite{Carzaniga2001}. Industry standards such as the Data Distribution Service (DDS) and MQTT have
further formalized deployment-time design choices---including topics, reliability guarantees,
durability, and brokered or brokerless architectures \cite{OMG2015DDS,OASIS2019MQTT}. While these
mechanisms enable the construction of complex modern cyber-physical, cloud, IoT, and robotics
architectures, they also make it difficult to reason about failure propagation using only direct
communication edges.

Consequently, research on pub-sub middleware dependability has traditionally emphasized fault
tolerance, reliable event dissemination, replication, and recovery strategies. More recent efforts
have employed graph-based dependency analyses to identify critical components within distributed
publish-subscribe environments. While these approaches give valuable guidance for design-time
reliability assessment, they do not offer the heterogeneous learning framework central to our
study. Rather than focusing on post-failure recovery, our work handles pre-deployment criticality
prediction: starting from an architectural model that enumerates applications, libraries, topics,
brokers, and Quality-of-Service (QoS) policies, we seek to estimate which application-level
components could have the greatest downstream impact if they were to fail. The view then requires a
graph that retains the semantics of the pub-sub model rather than collapsing dependencies into an
untyped network abstraction.

\subsection{Structural Criticality Analysis}

Classical graph analysis produces a rich toolkit for identifying important nodes and edges in
complex networks. Graph-based statistical metrics such as degree, closeness centrality, betweenness
centrality, articulation points, and PageRank-style scores are valued for their computational
efficiency along with interpretability \cite{Freeman1977,Brandes2001,BrinPage1998}. The field of
network science has further advanced our insight into system dependability by examining the effects
of node removals, chain failures, and interdependent network structures
\cite{Albert2000,Buldyrev2010}. In distributed systems, these measures commonly act as indicators of
components that act as critical communication bridges or whose removal could fragment the system,
making them useful reference points for pre-deployment architectural assessments.

However, relying exclusively on structural centrality provides only a limited view of failure risk
in publish-subscribe middleware. These systems are intentionally designed to decouple producers and
consumers via mechanisms such as topics, brokers, routing policies, and quality-of-service settings.
Consequently, a node that appears structurally central may not actually be semantically critical,
while a peripheral component could have outsized importance owing to its involvement in high-impact
topics or shared dependencies. In addition, conventional structural metrics commonly treat all edges
as identical, failing to distinguish how failures propagate through publication, subscription,
deployment, or logical dependency relations. While structural baselines supply valuable reference
points for evaluation, they lack sufficient expressiveness to comprehend the subtle, type-specific
semantics that are important for a complete and correct analysis of publish-subscribe middleware.

\subsection{Learning-Based Critical-Node Prediction}

Recent progress in graph learning has enabled more sophisticated strategies for detecting critical
nodes and uncovering structural patterns inside networks. Core models such as Graph Convolutional
Networks (GCN) and GraphSAGE learn node representations by aggregating information from neighboring
nodes \cite{Kipf2017,Hamilton2017}, while attention-based mechanisms dynamically modulate the
influence of neighbors during message passing \cite{Velickovic2018}. Using the mentioned
frameworks, FINDER~\cite{Fan2020} detects important entities in networked systems,
DrBC~\cite{Munikoti2022} predicts betweenness centrality using a graph neural network architecture,
and PowerGraph~\cite{Munikoti2024} applies graph learning to critical node analysis in power
systems. Thus, various studies show that learned graphs often outperform traditional and
human-crafted metrics, especially when higher-order structural features are at play.

However, most of these techniques are designed for homogeneous graphs or domain-based abstractions
in which all nodes and edges share a single semantic type. In contrast, publish-subscribe
middleware graphs are fundamentally heterogeneous: applications publish and subscribe to topics;
these topics are routed via brokers; libraries introduce software dependencies; and deployment
nodes impose locality constraints. Flattening such a structure into a homogeneous graph discards
critical information about how failures propagate and how the system behaves dynamically. The HGL
framework directly handles this drawback by learning relation-specific message functions over a
typed architectural model, consequently preserving the semantic richness and operational semantics
of middleware topologies.

\subsection{Heterogeneous Graph Neural Networks}

Heterogeneous graph neural networks (HGNNs) offer a well-founded approach to learning from graphs
with multiple node and edge types. Notable models include RGCN, which applies relation-specific
transformations in multi-relational graphs \cite{Schlichtkrull2018}; HAN, which implements layered
attention to aggregate information across both node and semantic levels \cite{Wang2019}; HGT, which
parameterizes attention according to node and edge type \cite{Hu2020}; and MAGNN, which utilizes
metapath-based aggregation to capture richer heterogeneous contexts \cite{Fu2020}. These models
have achieved strong performance in domains where the semantics of relations are central, such as
knowledge graphs, recommender systems, and information networks.

These advances are particularly pertinent for publish-subscribe middleware, which intrinsically
maps to a typed graph architecture: applications, topics, brokers, libraries, and deployment nodes
are linked by distinct relations, each with different consequences for failure propagation.
However, conventional message passing across densely connected or hub-dominated regions can lead to
over-smoothing, in which node representations become indistinguishable after multiple aggregation
steps \cite{Li2018}. This phenomenon is especially pronounced in pub-sub topologies with
high-fan-out topic and broker layers. Accordingly, our approach adopts heterogeneous learning
precisely in these contexts: HGL preserves the typed transport-level relationships in the native
architecture graph, while our prediction targets and evaluation remain at the application level.

\subsection{Positioning of This Work}

To summarize, the most relevant available approaches either examine pub-sub dependability at the
protocol or runtime level, offer interpretable but limited structural analyses, or employ
graph-learning techniques that overlook the typed semantics fundamental to publish-subscribe
middleware. HGL fills this gap via operating directly on graph models and learning how these
structural features influence failure propagation. Importantly, our contribution is not to propose
a new generic HGNN architecture, but to formulate, implement, and empirically evaluate a
middleware-specific application of heterogeneous graph learning for pre-deployment cascade impact
prediction.

% ===========================================================================
\section{Our Methodology}

This section details the methodology we developed to assess whether heterogeneous graph learning
(HGL) can reliably predict the impact of runtime failure cascades using pre-deployment graph-based
models of publish-subscribe middleware. The evaluation is anchored by the three research questions
enumerated below.

\textbf{RQ1.} Does graph learning outperform structural-centrality baselines such as betweenness,
articulation points, and QoS-weighted variants in predicting critical components within pub-sub
topologies?

\textbf{RQ2.} Does maintaining heterogeneous node and relation semantics confer benefits over a
homogeneous graph-learning baseline that treats all nodes and edges uniformly?

\textbf{RQ3.} Within the heterogeneous architecture, does including explicit QoS edge attributes
increase predictive performance relative to a QoS-masked heterogeneous model?

\subsubsection*{Formal Definitions}

For any given publish-subscribe system model, HGL assigns a \textbf{criticality score} $Q^*(v) \in
[0,1]$ to each component $v$ --- the model's predicted ranking of how damaging $v$'s failure would
be.

The ground-truth target, \textbf{cascade impact} $I^*(v) \in [0,1]$, is produced by an independent
discrete-event fault-injection engine (\texttt{FaultInjector}; implementation:
\texttt{saag/simulation/fault\_injector.py}), not by the learning models themselves. For each
candidate component $v$ and each seed $s \in \{42, 123, 456, 789, 2024\}$, the engine injects a
failure at $v$ and, for every subscriber application, computes a per-topic \textbf{feed-loss
fraction} --- the rate-weighted share of that topic's publishers that have failed, or the
equivalent broker-routing loss when the topic has no direct publisher among the failed set. A
\textbf{depth-damping} factor discounts loss propagated through indirect (multi-hop) paths relative
to direct ones. Each topic's feed-loss fraction is then scaled by a \textbf{topic QoS factor}:
$\times 1.2$ if the topic's reliability is \texttt{RELIABLE}, $\times 1.15$ if its transport
priority is \texttt{HIGH}/\texttt{CRITICAL}/\texttt{URGENT}, $\times 1.05$ if \texttt{MEDIUM},
clamped to $[0,1]$. $I^*(v)$ is the mean of these QoS-scaled feed-loss fractions across all of $v$'s
downstream subscribers and across seeds. A \textbf{propagation threshold} of $0.2$ governs whether a
subscriber's own loss is itself treated as a failure that can propagate further downstream. We call
this continuous construction \textbf{simulation softening}: a binary delivered/undelivered outcome
produces sparse, near-degenerate labels in decoupled pub-sub topologies (most components affect
almost nothing), so softening is what makes $I^*(v)$ informative enough to rank against.

We state this precisely, and in the paper body rather than only in the replication package, because
reviewers of the conference version of this work correctly flagged that a black-box description of
the simulator falls short of journal reproducibility standards. The anonymized/public replication
package (\S~3.4) remains available as a supplement --- the full failure-injection model,
per-scenario topology-generation scripts, and exact configuration --- but is not a substitute for
this definition.

To answer RQ1--RQ3, we use a controlled $2\times3$ factorial experimental design, methodically
varying both model design and QoS encoding, along with two non-learning structural baselines. The
evaluation comprises seven distinct scenarios and five independent random seeds, yielding
\textbf{210 evaluation cells}: 140 trained GNN models and 70 structural-baseline computations.
Table~\ref{tab:variants} summarizes the model variants and configurations analyzed in the study.

\begin{table}[htbp]
\centering
\caption{Model Variants and Baseline Configurations in the Controlled $2\times3$ Factorial Design}
\label{tab:variants}
\small
\begin{tabular}{@{}l p{2.6cm} p{2.8cm} p{4.5cm}@{}}
\toprule
Variant & Architecture & QoS Encoding & Description \\
\midrule
HGL-QoS & Heterogeneous GAT & 7-dimensional vector & Proposed method (full QoS encoding). \\
HGL & Heterogeneous GAT & masked & Proposed method (QoS-masked). \\
GL-QoS & Homogeneous GAT & scalar edge weight & Homogeneous baseline, scalar QoS weight. \\
GL & Homogeneous GAT & none & Homogeneous baseline, no QoS weight. \\
Topo-QoS & Structural centrality & QoS-weighted betweenness & Strongest structural baseline. \\
Topo-BL & Structural centrality & none & Unweighted betweenness \& articulation points. \\
\bottomrule
\end{tabular}
\normalsize
\end{table}

\subsection{Graph Representation}
\label{sec:graphrep}

Each deployment scenario is modeled as a heterogeneous directed graph:
\[
G = (V, E, \tau_V, \tau_E, w, \text{QoS})
\]
where $V$ is the set of architectural components and $E \subseteq V \times V$ is the set of typed
dependencies. The node- and edge-type vocabularies are:
\[
T_V=\{\text{Application}, \text{Library}, \text{Topic}, \text{Broker}, \text{Node}\}
\]
\begin{multline*}
T_E=\{\text{PUBLISHES\_TO}, \text{SUBSCRIBES\_TO}, \text{ROUTES}, \text{RUNS\_ON}, \\
\text{CONNECTS\_TO}, \text{USES}, \text{DEPENDS\_ON}\}
\end{multline*}

Of these seven types, six are \textbf{structural} --- imported directly from the topology JSON:
PUBLISHES\_TO, SUBSCRIBES\_TO, ROUTES, RUNS\_ON, CONNECTS\_TO, and USES. ROUTES (Broker
$\to$ Topic) captures broker routing responsibility; a broker's criticality is visible in the
heterogeneous graph only through this edge type, because broker failure impacts every topic it
routes. DEPENDS\_ON is \textbf{derived}: it is not imported from the topology but added by six
derivation rules in the pre-analysis step (see below); it is listed in $T_E$ because the HGL model
consumes it alongside the structural edges, yielding the seven relation types reflected in the
7-dimensional one-hot edge encoding.

The maps $\tau_V : V \to T_V$ and $\tau_E : E \to T_E$ assign these types to nodes and edges,
respectively. The scalar weight $w:E\to\mathbb{R}_+$ captures structural intensity derived from
publication frequency, message size, and subscriber fan-out. The QoS map
$\text{QoS}:E\to\mathcal{Q}$ assigns reliability, durability, and transport-priority attributes to
pub-sub edges where those attributes are meaningful.

The transport-level graph is transformed into a logical dependency graph by adding derived
\texttt{DEPENDS\_ON} edges. The direction convention is \textbf{dependent $\to$ dependency}: if
Application $A$ publishes to Topic $T$ and Application $B$ subscribes to $T$, then $B$ depends on
$A$, so a dependency $B \xrightarrow{\text{DEPENDS\_ON}} A$ is added. Similarly, if Application $A$
uses Library $L$, then $A$ depends on $L$, so the edge $A \xrightarrow{\text{DEPENDS\_ON}} L$ is
introduced.

We state this convention twice, deliberately, because it runs counter to the visual habit pub-sub
diagrams otherwise train: transport diagrams draw arrows for \emph{data flow}, publisher $\to$
subscriber, so a dependency arrow drawn the other way --- subscriber $\to$ publisher --- reads as
reversed to a middleware-literate audience, even though it is not. To make the contrast explicit:
\textbf{data flows $A \to B$; dependency points $B \to A$. A \texttt{DEPENDS\_ON} arrow is drawn
\emph{against} the direction of data flow}, because it encodes what fails if $A$ fails, not what $A$
sends. (\S~5.5 tests this convention directly: inverting \texttt{DEPENDS\_ON} does not merely
degrade a structural predictor built on the projection, it flips its correlation with ground-truth
cascade impact from strongly positive to strongly negative.) This operation maps pub-sub routing
into an application-level logical layer that serves as the input for the homogeneous baselines and
the structural metrics (e.g., betweenness, articulation points, bridge ratio); the heterogeneous HGL
model, by contrast, consumes the native typed architecture graph directly. Figure~\ref{fig:graphrep}
contrasts the two graph representations consumed by the learned models: HGL operates on the
heterogeneous architecture model, which preserves the full five-type vocabulary and typed relations
(panel a), whereas the homogeneous baseline GL operates on the lifted logical dependency subgraph,
in which the node and edge types are collapsed into a single graph view (panel b).

\begin{figure}[htbp]
\centering
\setlength{\unitlength}{0.85mm}
\begin{minipage}{0.5\textwidth}
\centering
\footnotesize (a) Architecture Model (HGL input)\\[2pt]
\begin{picture}(85,62)
% --- nodes ---
\put(2,50){\framebox(16,8){$a_1$}}
\put(2,30){\framebox(16,8){$a_2$}}
\put(2,10){\framebox(16,8){$a_3$}}
\put(46,32){\oval(20,12)}
\put(44,30){\makebox(4,4){$t$}}
\put(36,0){\framebox(20,8){$\ell$}}
\put(64,50){\framebox(20,8){}}
\put(65,51){\framebox(18,6){$b$}}
\put(66,10){\framebox(16,8){$n$}}
% --- edges: a1 -> t (pub) ---
\put(18,54){\line(1,0){28}}
\put(46,54){\vector(0,-1){16}}
\put(30,56){\makebox(0,0){\tiny pub}}
% --- edges: a2 -> t (sub) ---
\put(18,34){\vector(1,0){18}}
\put(27,36){\makebox(0,0){\tiny sub}}
% --- edges: a3 -> t (sub) ---
\put(18,17){\line(1,0){28}}
\put(46,17){\vector(0,1){9}}
\put(30,19){\makebox(0,0){\tiny sub}}
% --- edges: b -> t (routes) ---
\put(70,50){\line(0,-1){18}}
\put(70,32){\vector(-1,0){14}}
\put(74,40){\makebox(0,0){\tiny routes}}
% --- edges: b -> n (runs) ---
\put(78,50){\vector(0,-1){32}}
\put(82,34){\makebox(0,0){\tiny runs}}
% --- edges: a2 -> l (uses) ---
\put(18,32){\line(1,0){4}}
\put(22,32){\line(0,-1){28}}
\put(22,4){\vector(1,0){14}}
\put(24,7){\makebox(0,0){\tiny uses}}
% --- edges: a3 -> l (uses) ---
\put(18,11){\line(1,0){10}}
\put(28,11){\line(0,-1){5}}
\put(28,6){\vector(1,0){8}}
\end{picture}
\end{minipage}
\hfill
\begin{minipage}{0.45\textwidth}
\centering
\footnotesize (b) Logical Dependency Subgraph (GL input)\\[2pt]
\begin{picture}(75,62)
% --- nodes (uniform circles, untyped), on a square grid: a2 top-left, a1 top-right,
% --- a3 bottom-left, l bottom-right, so a3->a1 and a2->l are exact 45-degree diagonals ---
\put(50,50){\circle{16}}
\put(50,50){\makebox(0,0){$a_1$}}
\put(10,50){\circle{16}}
\put(10,50){\makebox(0,0){$a_2$}}
\put(10,10){\circle{16}}
\put(10,10){\makebox(0,0){$a_3$}}
\put(50,10){\circle{16}}
\put(50,10){\makebox(0,0){$\ell$}}
% --- edges: a2 -> a1 (depends on) ---
\put(18,50){\vector(1,0){24}}
\put(30,54){\makebox(0,0){\tiny depends on}}
% --- edges: a3 -> a1 (depends on), 45-degree diagonal ---
\put(16,16){\vector(1,1){28}}
% --- edges: a2 -> l (depends on), 45-degree diagonal ---
\put(16,44){\vector(1,-1){28}}
\put(46,26){\makebox(0,0){\tiny depends on}}
% --- edges: a3 -> l (depends on) ---
\put(18,10){\vector(1,0){24}}
\end{picture}
\end{minipage}
\caption{The two graph representations consumed by the learned models. (a) HGL ingests the
\emph{heterogeneous} architecture model, which preserves the full five-type components---Applications
($a_i$, plain boxes), a Library ($\ell$, box), a Topic ($t$, ellipse), a Broker ($b$, double box),
and a compute Node ($n$, box)---and the typed relationships (pub, sub, routes, uses, runs), drawn in
the direction data actually flows (publisher $\to$ topic $\to$ subscriber; broker $\to$ topic;
broker $\to$ node). (b) The homogeneous baseline GL ingests the lifted logical dependency subgraph,
in which the transport tier is collapsed into derived \texttt{DEPENDS\_ON} edges and \emph{all node
and edge types are flattened into a single graph view} (uniform circles, untyped edges, labeled
``depends on'' to visually distinguish them from panel (a)'s typed pub/sub/routes/runs/uses labels).
Edges point from \emph{dependent} to \emph{dependency} --- the reverse of data flow: $a_2$ and $a_3$
both depend on publisher $a_1$ (subscriber $\to$ publisher arrows, against the pub $\to$ topic $\to$
sub data path), and both also depend on the shared library $\ell$ (user $\to$ library arrows). The
contrast between the two panels isolates the central design question: whether retaining node and
relation \emph{type} (a) yields better cascade prediction than reasoning over connectivity alone
(b). In this example, $a_1$'s role as the upstream publisher on which $a_2$ and $a_3$ depend, and
the shared library $\ell$ concentrating a multi-application blast radius, are signals that are
explicit in (a) but indistinguishable from ordinary edges in (b).}
\label{fig:graphrep}
\end{figure}

Each edge is encoded as a 16-dimensional feature vector: this includes a scalar structural weight,
a normalized path-count feature, a 7-dimensional one-hot encoding for edge type, and 7 QoS-derived
features (covering reliability, durability, transport priority, deadline, and lifespan). QoS
features are set to zero for non-pub/sub edges. The \texttt{HGL} variant masks QoS features on
pub-sub edges, while \texttt{HGL-QoS} exposes them, thereby isolating the value added by explicit
QoS attributes.

\subsection{Model Variants}

HGL is realized as a heterogeneous graph attention network, where each relation type is assigned a
dedicated message function. Message passing operates over the complete typed architecture: all
five node types (Applications, Libraries, Topics, Brokers, and Nodes) contribute to the aggregation
process, and each transport relation (publication, subscription, library usage, broker
connectivity, and host deployment) is uniquely parameterized. Notably, publish-subscribe edges
carry Quality-of-Service (QoS) attributes where applicable. By maintaining those distinctions
throughout aggregation, the model understands propagation patterns specific to each architectural
layer, rather than relying on a single uniform transformation. In contrast, the homogeneous
baselines (\texttt{GL} and \texttt{GL-QoS}) collapse all node and edge types into a unified graph
view and apply analogous attention models. These baselines are natural adaptations of standard
homogeneous GNNs (such as FINDER, DrBC, and PowerGraph) to middleware contexts, but they lack the
capacity for relation-specific message passing.

The factorial experimental design enables three focused comparisons: (1) \texttt{GL} versus
\texttt{HGL}, which isolates the impact of architectural heterogeneity while keeping QoS features
masked; (2) \texttt{HGL} versus \texttt{HGL-QoS}, which examines the effect of explicit QoS
encoding while holding the overall architecture constant; and (3) \texttt{Topo-BL} versus
\texttt{Topo-QoS}, which assesses the degree to which non-learning structural metrics can recover
QoS information. Together, these comparisons help disentangle the particular contributions of typed
graph learning and QoS attribute inclusion.

\subsection{Training and Validation}

Our study spans seven distinct deployment scenarios, encompassing domains such as autonomous
vehicles, high-frequency financial trading, clinical healthcare integration, centralized
hub-and-spoke enterprise systems, distributed IoT smart-city telemetry, cloud-native
microservices, and large-scale enterprise pub-sub architectures. Each scenario is constructed
using a synthetically generated topology, designed to reflect realistic counts of applications,
libraries, topics, brokers, and nodes. For example, the Microservices scenario comprises 32
applications, 6 libraries, 25 topics, 4 brokers, and 8 compute nodes.

Model training is performed using the PyTorch Geometric HeteroGAT implementation, with 4 attention
heads per relation, a hidden-layer dimension of 64, and 300 epochs per experimental cell. For each
random seed, nodes are partitioned into training and validation sets, with stratification by node
type to ensure uniform representation. Results are averaged over five independent seeds, and
uncertainty is quantified using bootstrap-derived 95\% confidence intervals ($B=2000$ resamples).
Identification metrics (F1, precision, recall, and Top-$K$ overlap) are computed using
rank-matched binarization: the top-$K$ predicted components are labeled as critical, where $K$ is
set to the number of ground-truth critical components ($I^*(v)>0.5$). Statistical significance
between HGL and each comparator is evaluated using paired Wilcoxon signed-rank tests across the
five seeds for each scenario.

We intentionally frame our empirical findings as relative architectural comparisons. The results
prove that HGL consistently improves the detection of critical components, outperforming both
structural and homogeneous graph-learning baselines. The $2\times3$ factorial design further shows
that these gains are driven primarily by the use of heterogeneous, typed message passing, rather
than the explicit inclusion of QoS attributes. It is important to clarify that we do not assert
that the absolute values of $\rho$ or F1 will generalize to any arbitrary deployed pub-sub system.
Both the predicted scores $Q^*(v)$ and the ground-truth labels $I^*(v)$ are determined within the
confines of the same experimental framework, and their absolute alignment is inherently limited by
the accuracy and fidelity of the simulator.

\subsection{Cascade Impact Simulation}

The ground-truth impact $I^*(v)$ is produced by \texttt{FaultInjector}, a discrete-event cascade
simulator (formally defined in \S~3) that operates on the same publish-subscribe system model
consumed by the learning algorithms, but through an independent process: for each candidate
component $v$, the simulator injects a failure at $v$ and propagates the resulting disruption along
publication and subscription paths over a fixed simulation horizon, while the learning models never
observe the simulator's internal state. This is the sole engine used to produce every $I^*(v)$
label reported in this paper --- a separate AHP-weighted composite-impact metric exists elsewhere
in the codebase for CI/CD reporting purposes but is not used for, and should not be conflated with,
the labels evaluated here. Impact is quantified as the degradation of communication at subscriber
applications relative to the fault-free baseline. Rather than recording a binary
delivered/undelivered outcome---which yields sparse, near-degenerate labels in decoupled
topologies---we apply a \emph{simulation-softening} procedure that produces a continuous target
$I^*(v)\in[0,1]$ from rate-weighted fractions of failed publishers combined with topic-level QoS
factors (reliability and transport priority).

To make these settings fully reproducible, the complete simulator configuration (including the
failure-injection model, propagation rules, QoS-factor weighting, simulation horizon, and the
per-scenario topology-generation scripts) is provided in an anonymized replication
package.\footnote{\url{https://anonymous.4open.science/r/software-as-a-graph-A622/reproduce/README.md}}
The package covers the full evaluation: the main 210-cell factorial harness behind Tables 1 and
4--7, and standalone scripts for the reversed-projection ablation and hardening-budget analysis
(\S~5.5--5.6, Tables 8--9).

% ===========================================================================
\section{Evaluation}

The evaluation suite is structured to determine whether a pre-deployment graph model can
meaningfully guide two critical architectural decisions: (1) ranking system components by their
anticipated cascade impact and (2) identifying which components have to be prioritized for
hardening interventions. To present a holistic assessment of model results, we report metrics
spanning ranking, identification, and regression. Furthermore, we implement a controlled masking
protocol to clearly separate the effects of architectural structure from those attributable to
Quality-of-Service (QoS) features.

\paragraph{Evaluation Protocol.} Each of the seven scenarios (\S~4.6) is trained and evaluated
independently, in two regimes. \emph{In-distribution}: for each scenario and each of five seeds
$\{42, 123, 456, 789, 2024\}$, nodes are split into training and validation sets stratified by node
type, the model variant (Table~\ref{tab:variants}) is trained on that scenario's graph alone, and
metrics (\S~4.1--4.3) are computed against that same scenario's held-out validation split.
\emph{Leave-One-Scenario-Out (LOSO, \S~5.4)}: for each held-out scenario, the model trains on the
graphs of the remaining six and is evaluated, without any further tuning, on the untouched held-out
scenario --- the out-of-distribution regime a model faces against a genuinely new deployment. In
both regimes, input--label independence is structural: the GNN's edge and node features (\S~3.1)
are computed from the same topology JSON that FaultInjector consumes, but never from FaultInjector's
own internal simulation state (visited sets, propagation traces, or intermediate feed-loss values)
--- only the final $I^*(v)$ label crosses from simulator to evaluation, never into training
features.

\subsection{Ranking Metrics}

Our main ranking metric is the Spearman rank correlation ($\rho$), which quantifies how well the
ordering of predicted criticality scores $Q^*(v)$ corresponds to the simulator-derived impact
scores $I^*(v)$ for each scenario and random seed. This approach is particularly appropriate for
pre-deployment analysis, where the primary concern is ordering components by relative criticality
rather than determining their exact failure-impact magnitudes. We summarize performance by
reporting the average $\rho$ across seeds and provide bootstrap-derived 95\% confidence intervals
to reflect statistical variability.

To provide additional perspective, we report \textbf{NDCG@10} (Normalized Discounted Cumulative
Gain at rank 10) as a complementary top-of-list metric. Whereas Spearman's correlation assesses
alignment across the entire ranking, NDCG@10 specifically measures how effectively the model
identifies the most impactful components at the top of the list, a key consideration when
operational constraints restrict the number of components that can be reviewed or hardened.

\subsection{Critical-Component Identification}

Identification metrics assess a model's ability to correctly identify the components that require
prioritized hardening interventions. We define the ground-truth critical set as:
\[
G = \{v : I^*(v) > 0.5\}
\]
with its size denoted $K=|G|$. For each model, the predicted set $P$ is formed by selecting the
top-$K$ components according to their predicted criticality scores $Q^*(v)$. This
\textbf{rank-matched binarization} strategy decouples the ranking task from score calibration,
since models may output scores on different, non-comparable scales. Using a fixed threshold such as
$Q^*(v)>0.5$ could conflate calibration artifacts with true predictive performance.

With rank-matched binarization, the sizes of $P$ and $G$ are always equal ($|P|=|G|=K$). When
$0<K<|V|$, the standard identification metrics---precision, recall, and F1-score---all reduce to
the same value:
\[
\text{Precision}=\text{Recall}=\text{F1}=\frac{|P\cap G|}{K}
\]

Accordingly, F1 is reported as the primary identification metric. Edge cases where $K=0$ or
$K=|V|$ do not occur in our scenarios, but the evaluation framework is designed to detect and flag
such situations if they arise.

Accuracy is reported only as a secondary metric for completeness and comparability. It is not
emphasized in our main results because it is highly sensitive to the prevalence of critical
components: when only a few components are classified as critical, a trivial model could achieve
deceptively high accuracy by predicting all components as non-critical.

\subsection{Regression Error}

Regression metrics quantify the discrepancy between projected and actual component scores. In
particular, we report \textbf{RMSE} (Root Mean Squared Error), which disproportionately penalizes
larger errors, and \textbf{MAE} (Mean Absolute Error), which represents the average magnitude of
prediction error. While these metrics provide important insights into model calibration and
accuracy, we treat them as secondary to ranking and identification measures. This prioritization
stems from the fact that the absolute scale of $I^*(v)$ is a function of simulator-specific
settings; accordingly, our empirical focus is on relative performance in ranking and
identification, rather than on absolute regression error.

\subsection{Statistical Aggregation}

To ensure robust evaluation, each model variant is assessed with five independent random seeds per
scenario. Metrics are averaged within each experimental cell, and uncertainty is captured through
bootstrap-derived 95\% confidence intervals. For statistical comparisons between HGL and baseline
approaches, we apply paired Wilcoxon signed-rank tests at the seed level. Since all scenarios,
seeds, graph data, and simulation targets remain fixed across variants, these pairwise comparisons
provide direct and controlled estimates of the impact of architectural and feature-encoding
decisions on model results.

\subsection{Ablation Protocol}

Our ablation protocol is structured to separate the particular contributions of heterogeneous graph
architecture and explicit QoS attributes. The no-QoS variants---\textbf{GL} (\texttt{gl}) and
\textbf{HGL} (\texttt{hgl})---are trained on graph features in which QoS information has been
masked before conversion to PyTorch Geometric format. Conversely, the QoS-aware
variants---\textbf{GL-QoS} (\texttt{gl\_qos}) and \textbf{HGL-QoS} (\texttt{hgl\_qos})---are trained
using the complete QoS-annotated graph. This experimental setup allows us to isolate the effect of
architectural heterogeneity (by comparing \texttt{GL} to \texttt{HGL}) and the incremental value of
including explicit QoS information (by comparing \texttt{HGL} to \texttt{HGL-QoS}).

\subsection{Domain Scenario Taxonomy and Dataset Scale}

To comprehensively evaluate the predictive performance and robustness of HGL across a variety of
architectural patterns, we developed an evaluation suite comprising seven distinct scenarios. These
scenarios span six major topology classes and domain verticals, including fan-out-dominated
systems, dense pub-sub networks, single-point-of-failure (SPOF) architectures, sparse,
low-coupling deployments, safety-critical, real-time environments, and over-provisioned, redundant
broker topologies. Each benchmark is synthesized using parameterized configurations informed by
real-world middleware deployment constraints and domain models, ensuring representative diversity
in topology density, Quality-of-Service (QoS) settings, and structural complexity.

The evaluation suite is organized according to a multi-tier, scale-preset taxonomy, ranging from
ultra-compact smoke tests to a hyper-scale enterprise platform with up to 300 applications and over
one hundred communication channels. Table~\ref{tab:scale} summarizes the core structural
characteristics and tier allocations for these scenario profiles.

These seven scenarios are derived from the standard scale presets while maintaining unique
domain-specific wiring layouts. (This evaluation suite deliberately excludes the ``Tiny Regression''
CI smoke-test fixture used elsewhere in the pipeline for deterministic anti-pattern detection
regression testing, since it carries no domain-representative topology and would not contribute
meaningful predictive signal.)

\begin{itemize}
\item \textbf{Autonomous Vehicle (AV) System (Scenario 01):} Built on the \texttt{medium} scale,
  representing a high fan-out ROS2 autonomous vehicle system with sensor streaming topologies
  heavily bound to \texttt{RELIABLE} and \texttt{TRANSIENT\_LOCAL} QoS profiles.
\item \textbf{IoT Smart City (Scenario 02):} Configured at the \texttt{large} scale, characterizing
  massive, distributed endpoint matrices that generate high-frequency message streams bound to
  high-loss \texttt{VOLATILE} and \texttt{BEST\_EFFORT} transport contracts.
\item \textbf{Financial Trading (Scenario 03):} Built on a dense \texttt{medium} network schema
  with critical message flows executing under strict high-priority \texttt{PERSISTENT} persistence
  rules.
\item \textbf{Healthcare Systems (Scenario 04):} Operates as a dense \texttt{medium} deployment
  reflecting clinical health integration architectures, characterized by centralized patient
  monitoring fan-outs and long-lived durability contracts.
\item \textbf{Hub-and-Spoke (Scenario 05):} A deliberate structural anti-pattern framework built at
  the \texttt{medium} count boundary, constraining system routing through exactly two brokers to
  evaluate the engine's capability to flag single points of failure (SPOF).
\item \textbf{Microservices (Scenario 06):} Represents a cloud-native service mesh deployed as a
  sparse topology with low architectural coupling to benchmark model correctness and guard against
  classification threshold over-flagging.
\item \textbf{Enterprise (Scenario 07):} A complex \texttt{xlarge} scale environment (comprising
  300 distinct applications) serving as our primary performance and scalability bottleneck check.
\end{itemize}

\begin{table}[htbp]
\centering
\caption{Dataset Scale Presets and Structural Tier Distributions}
\label{tab:scale}
\begin{tabular}{@{}lrrrrrl@{}}
\toprule
Scale Preset & Apps & Libs & Topics & Brokers & Nodes & Target Scenario \\
\midrule
\texttt{small} & 15 & 5 & 10 & 2 & 4 & General low-scale baseline, smoke tests \\
\texttt{medium} & 50 & 10 & 30 & 3 & 8 & Scenarios 01, 03, 04, 05, 06 \\
\texttt{large} & 150 & 30 & 100 & 6 & 20 & Scenario 02 (IoT Smart City) \\
\texttt{xlarge} & 300 & 50 & 120 & 10 & 40 & Scenario 07 (Enterprise) \\
\bottomrule
\end{tabular}
\end{table}

Table~\ref{tab:scale} reports the \emph{preset} targets each scenario is drawn from; the generator
introduces per-scenario variance on top of these targets to reflect domain-specific wiring.
Table~\ref{tab:scenariochar} reports the \emph{as-generated} structural counts actually evaluated,
together with the deployment pattern each scenario is meant to characterize, so that a reader can
judge representativeness without cross-referencing the replication package.

\begin{table}[htbp]
\centering
\small
\caption{Scenario Characterization --- As-Generated Structure and QoS Mix}
\label{tab:scenariochar}
\begin{tabular}{@{}lrrll@{}}
\toprule
Scenario & $|V|$ (App/Lib/Topic/Broker/Node) & $|E|$ (total) & Edge mix (pub/sub/routes/uses/runs/conn) & Topic QoS mix \\
\midrule
AV System & 152 (80/20/40/4/8) & 797 & 146/313/53/142/129/14 & 40 RELIABLE \\
IoT Smart City & 326 (200/10/80/6/30) & 1322 & 374/309/103/82/148 & 35 REL., 45 BEST\_EFFORT \\
Financial Trading & 124 (60/18/35/5/6) & 580 & 144/187/47/103/94/5 & 29 REL., 6 BEST\_EFFORT \\
Healthcare & 98 (50/12/25/3/8) & 400 & 95/111/31/73/82/8 & 19 REL., 6 BEST\_EFFORT \\
Hub-and-Spoke & 139 (70/25/30/2/12) & 797 & 140/310/42/182/109/14 & 30 RELIABLE \\
Microservices & 186 (90/30/45/6/15) & 680 & 149/221/58/82/138/32 & 45 RELIABLE \\
Enterprise & 520 (300/50/120/10/40) & 3245 & 769/1206/151/437/449/233 & 120 RELIABLE \\
\bottomrule
\end{tabular}
\normalsize
\end{table}

The real deployment pattern each scenario represents is described in the itemized list above;
in brief, and referencing the structural counts in Table~\ref{tab:scenariochar}: AV System is a
ROS2-style sensor fan-out with many subscribers per topic (sub:pub~$\approx$~2.1) and
all-RELIABLE QoS reflecting safety-critical sensor streaming; IoT Smart City is a large,
distributed endpoint matrix with majority-BEST\_EFFORT QoS modeling high-loss telemetry links;
Financial Trading is a dense trading mesh with the highest \texttt{uses}-edge density among
medium-scale scenarios (103), reflecting heavy shared-library reliance typical of trading stacks;
Healthcare is the smallest scenario by $|V|$, a centralized patient-monitoring fan-out with
long-lived (\texttt{PERSISTENT}) durability contracts; Hub-and-Spoke is a deliberate SPOF
anti-pattern routing 30 topics through only 2 brokers, with the highest overall \texttt{uses}-edge
density (182) concentrating blast radius in shared libraries; Microservices is a sparse,
low-coupling service mesh with the lowest sub:pub ratio (1.48), used to check against
over-flagging; and Enterprise is the hyper-scale platform (largest $|V|$ and $|E|$) serving as the
primary scalability/performance bottleneck check.

Generation parameters (scale preset, seed, per-scenario domain wiring script) and the exact
topology JSON for each row are provided in the replication package (\S~3.4), under
\texttt{data/scenarios/} as \texttt{av\_system.json}, \texttt{iot\_smart\_city\_system.json},
\texttt{financial\_trading\_system.json}, \texttt{healthcare\_system.json},
\texttt{hub\_and\_spoke\_system.json}, \texttt{microservices\_system.json}, and
\texttt{enterprise\_system.json}.

% ===========================================================================
\section{Experimental Results}

The evaluation yields three principal findings, and the third qualifies the other two. First,
within a trained scenario's distribution, HGL effectively balances the tasks of ranking components
by potential cascade impact and accurately identifying those that should be prioritized for
hardening, with the most notable improvements in the identification task, where heterogeneous node
and relation types yield substantially better results than homogeneous graph learning. Second,
explicit QoS edge attributes do not improve in-distribution performance, suggesting the typed graph
structure already captures most of the routing information QoS features would add. Third --- and
this is the finding that qualifies the first two --- Leave-One-Scenario-Out evaluation shows that
this in-distribution advantage does not transfer to genuinely unseen scenarios: both HGL and
HGL-QoS fall to negative mean rank correlation out-of-distribution, underperforming even the
homogeneous baselines. We report this as it came out, not adjusted to fit the in-distribution
story, because a generalization gap this large is exactly the kind of finding an
in-distribution-only evaluation would miss.

From a quantitative standpoint, HGL attains the highest mean in-distribution ranking correlation
among the learned variants (Spearman $\rho=0.411$), behind only the QoS-weighted structural
baseline (Topo-QoS, $\rho=0.612$) in absolute terms, but clearly ahead of the homogeneous GL
($\rho=0.186$) and GL-QoS ($\rho=0.219$). For the identification task, HGL achieves a mean F1-score
of 0.688 in identifying critical components --- an improvement of $\Delta\text{F1}=+0.092$ over the
homogeneous GL baseline. In Leave-One-Scenario-Out validation, however, both heterogeneous variants
underperform both homogeneous baselines: HGL-QoS's mean out-of-distribution $\rho=-0.058$ and
HGL's $\rho=-0.114$, against GL's $\rho=0.021$ and GL-QoS's $\rho=0.002$ (\S~5.4).

These findings suggest that HGL's in-distribution advantage lies in its superior preservation of
middleware semantics while training and test data share a scenario's topology. While homogeneous
GNNs capture only raw connectivity, HGL maintains distinctions among all five node types
(Applications, Libraries, Topics, Brokers, and Nodes) and the specific relations connecting
them---such as publication, subscription, library usage, broker connectivity, and host
deployment. By retaining these distinctions throughout message passing, HGL delivers notably
stronger in-distribution performance in identifying components that should be prioritized for
hardening. What that richer representation does not do, on this evidence, is transfer those
distinctions to a scenario the model has never trained on --- \S~5.4 examines why.

Ablation analyses provide further insight into the in-distribution improvements. When QoS features
are masked for both models, switching from \texttt{GL} to \texttt{HGL} results in a mean ranking
increase of $\Delta\rho=+0.225$ and a mean identification gain of $\Delta\text{F1}=+0.092$. In
contrast, adding explicit QoS edge attributes to the heterogeneous model produces only marginal
differences and, on average, slightly reduces performance (\texttt{HGL-QoS} vs.\ \texttt{HGL}:
$\Delta\rho=-0.035$, $\Delta\text{F1}=-0.004$). This pattern indicates that typed heterogeneous
message passing is the primary architectural factor driving in-distribution performance, while
explicit QoS features contribute little there --- and, as \S~5.4 shows, do not rescue
out-of-distribution performance either.

\subsection{Ranking Performance}

The following table summarizes the correlation between global rankings across all scenarios and
variants.

\begin{table}[htbp]
\centering
\caption{Global Ranking Performance (Spearman $\rho$) Across Scenarios}
\label{tab:ranking}
\begin{tabular}{@{}lrrrrrrr@{}}
\toprule
Scenario & Topo-BL & Topo-QoS & GL & GL-QoS & HGL & HGL-QoS & $\Delta\rho$ (QoS) \\
\midrule
AV System & 0.485 & 0.691 & 0.352 & 0.318 & \textbf{0.496} & 0.455 & $-0.041$ \\
Enterprise & 0.361 & 0.731 & 0.376 & 0.421 & 0.832 & \textbf{0.863} & $+0.031$ \\
Financial Trading & 0.056 & 0.514 & 0.167 & 0.073 & \textbf{0.306} & 0.270 & $-0.036$ \\
Healthcare & 0.188 & 0.702 & 0.063 & 0.003 & \textbf{0.069} & $-0.116$ & $-0.185$ \\
Hub-and-Spoke & 0.235 & 0.813 & $-0.045$ & 0.172 & \textbf{0.249} & 0.144 & $-0.105$ \\
IoT Smart City & 0.065 & 0.255 & 0.382 & 0.504 & 0.695 & \textbf{0.745} & $+0.050$ \\
Microservices & 0.239 & 0.578 & 0.008 & 0.040 & 0.234 & \textbf{0.276} & $+0.042$ \\
\midrule
Mean & 0.233 & 0.612 & 0.186 & 0.219 & \textbf{0.411} & 0.377 & $-0.035$ \\
\bottomrule
\end{tabular}
\end{table}

Within a trained scenario's distribution, HGL and HGL-QoS consistently outrank the homogeneous
GL/GL-QoS baselines, though not the QoS-weighted structural baseline Topo-QoS, which remains the
strongest single predictor in five of seven scenarios (Table~\ref{tab:ranking}) --- a point we
return to in \S~6. In the Enterprise scenario, HGL-QoS attains its best correlation of
\textbf{0.863}, a \textbf{+0.487} improvement over GL (0.376). In IoT Smart City, HGL-QoS reaches
\textbf{0.745}, versus GL's 0.382 (\textbf{+0.363}). The QoS-encoding effect ($\Delta\rho$,
rightmost column) is inconsistent in sign and, in Healthcare, sizeable and negative ($-0.185$);
\S~5.3 examines this ablation systematically rather than scenario-by-scenario. Across all seven
scenarios, the heterogeneous GAT models outperform their homogeneous counterparts, but the margin
over Topo-QoS is narrow to negative --- unlike the conference submission's reported numbers, this
revision's HGL variants do not consistently exceed the QoS-weighted structural baseline
in-distribution.

\subsection{Identification Metrics}

The following table provides a breakdown of binary classification performance for critical
component identification.

\begin{longtable}{@{}llrrrrr@{}}
\caption{Identification and Regression Metrics Across the 7 Scenarios}
\label{tab:identification} \\
\toprule
Scenario & Variant & F1 & Accuracy & RMSE & MAE & NDCG@10 \\
\midrule
\endfirsthead
\multicolumn{7}{l}{\textit{Table~\ref{tab:identification} continued from previous page}} \\
\toprule
Scenario & Variant & F1 & Accuracy & RMSE & MAE & NDCG@10 \\
\midrule
\endhead
\midrule
\multicolumn{7}{r}{\textit{continued on next page}} \\
\endfoot
\bottomrule
\endlastfoot
AV System & Topo-BL & 0.775 & 0.775 & 0.581 & 0.549 & 0.666 \\
 & Topo-QoS & 0.825 & 0.825 & 0.581 & 0.550 & 0.660 \\
 & GL & 0.633 & 0.600 & 0.449 & 0.419 & 0.912 \\
 & GL-QoS & 0.600 & 0.564 & 0.473 & 0.438 & 0.901 \\
 & HGL & 0.767 & 0.745 & 0.184 & 0.158 & 0.936 \\
 & HGL-QoS & 0.767 & 0.745 & 0.196 & 0.163 & 0.929 \\
\addlinespace
Enterprise & Topo-BL & 0.649 & 0.640 & 0.590 & 0.567 & 0.887 \\
 & Topo-QoS & 0.812 & 0.807 & 0.590 & 0.567 & 0.899 \\
 & GL & 0.660 & 0.651 & 0.359 & 0.325 & 0.730 \\
 & GL-QoS & 0.670 & 0.662 & 0.296 & 0.253 & 0.783 \\
 & HGL & 0.840 & 0.836 & 0.147 & 0.125 & 0.942 \\
 & HGL-QoS & 0.870 & 0.867 & 0.154 & 0.127 & 0.956 \\
\addlinespace
Financial Trading & Topo-BL & 0.581 & 0.567 & 0.621 & 0.590 & 0.539 \\
 & Topo-QoS & 0.677 & 0.667 & 0.622 & 0.591 & 0.644 \\
 & GL & 0.560 & 0.511 & 0.447 & 0.410 & 0.907 \\
 & GL-QoS & 0.520 & 0.467 & 0.497 & 0.453 & 0.895 \\
 & HGL & 0.640 & 0.600 & 0.214 & 0.191 & 0.937 \\
 & HGL-QoS & 0.600 & 0.556 & 0.221 & 0.190 & 0.932 \\
\addlinespace
Healthcare & Topo-BL & 0.577 & 0.560 & 0.585 & 0.549 & 0.642 \\
 & Topo-QoS & 0.808 & 0.800 & 0.586 & 0.552 & 0.765 \\
 & GL & 0.627 & 0.543 & 0.467 & 0.407 & 0.879 \\
 & GL-QoS & 0.677 & 0.600 & 0.399 & 0.341 & 0.878 \\
 & HGL & 0.600 & 0.543 & 0.216 & 0.172 & 0.865 \\
 & HGL-QoS & 0.600 & 0.543 & 0.225 & 0.189 & 0.852 \\
\addlinespace
Hub-and-Spoke & Topo-BL & 0.605 & 0.571 & 0.563 & 0.527 & 0.739 \\
 & Topo-QoS & 0.842 & 0.829 & 0.564 & 0.528 & 0.701 \\
 & GL & 0.500 & 0.480 & 0.429 & 0.379 & 0.872 \\
 & GL-QoS & 0.607 & 0.600 & 0.428 & 0.378 & 0.882 \\
 & HGL & 0.520 & 0.520 & 0.209 & 0.170 & 0.921 \\
 & HGL-QoS & 0.520 & 0.520 & 0.210 & 0.169 & 0.900 \\
\addlinespace
IoT Smart City & Topo-BL & 0.480 & 0.480 & 0.588 & 0.565 & 0.528 \\
 & Topo-QoS & 0.510 & 0.510 & 0.589 & 0.565 & 0.615 \\
 & GL & 0.657 & 0.644 & 0.499 & 0.471 & 0.866 \\
 & GL-QoS & 0.700 & 0.689 & 0.571 & 0.547 & 0.880 \\
 & HGL & 0.814 & 0.807 & 0.141 & 0.115 & 0.949 \\
 & HGL-QoS & 0.829 & 0.822 & 0.126 & 0.100 & 0.956 \\
\addlinespace
Microservices & Topo-BL & 0.660 & 0.644 & 0.614 & 0.589 & 0.756 \\
 & Topo-QoS & 0.787 & 0.778 & 0.614 & 0.589 & 0.822 \\
 & GL & 0.533 & 0.491 & 0.389 & 0.351 & 0.841 \\
 & GL-QoS & 0.467 & 0.418 & 0.500 & 0.470 & 0.900 \\
 & HGL & 0.633 & 0.600 & 0.192 & 0.150 & 0.916 \\
 & HGL-QoS & 0.600 & 0.564 & 0.182 & 0.150 & 0.912 \\
\end{longtable}

The identification results are more mixed than the ranking correlations alone suggest. The
heterogeneous graph-learning models (HGL and HGL-QoS) achieve the highest F1 in four of seven
scenarios --- Enterprise (\textbf{0.840}/0.870), IoT Smart City (0.814/\textbf{0.829}), Financial
Trading (\textbf{0.640}), and, tied with Topo-BL, AV System (\textbf{0.767}) --- but the
QoS-weighted structural baseline Topo-QoS wins outright in Healthcare (0.808), Hub-and-Spoke
(0.842), and Microservices (0.787), each by a substantial margin over every learned variant. On
average, HGL achieves a mean F1 score of \textbf{0.688}, outperforming GL (\textbf{0.596}) by
$\Delta\text{F1} = +0.092$ --- a real but considerably smaller margin than ranking correlation
alone would suggest, and smaller than Topo-QoS's own mean F1 of 0.752. These results indicate that
heterogeneous message passing improves on homogeneous GNN baselines in-distribution, but does not,
on this evidence, establish HGL as uniformly superior to a QoS-weighted structural heuristic within
the trained distribution --- that comparison reverses out-of-distribution, where the structural
baseline's determinism becomes an advantage rather than a limitation (\S~5.4).

The practical implication is that in-distribution superiority is scenario-dependent rather than
uniform: for a system architect validating against historical, similar-topology deployments, HGL
and HGL-QoS are the stronger choice in roughly half of the evaluated deployment patterns, while the
cheaper, deterministic Topo-QoS baseline remains competitive or superior in the other half,
particularly in the SPOF-oriented Hub-and-Spoke and dense-fan-out Healthcare and Microservices
scenarios.

\subsection{Ablation Analysis}

The $2\times3$ factorial design (architecture $\times$ QoS encoding), together with the two
structural baselines, enables four controlled comparisons. These comparisons break down the
central finding reported above --- that HGL is the strongest learned variant across both ranking
and identification tasks in-distribution --- into its key architectural and QoS-encoding
contributions. In each comparison, one factor is held constant while the other is varied; the
resulting $\Delta\rho$ and $\Delta\text{F1}$ values, calculated from scenario-level means, are
presented in Table~\ref{tab:ablation}.

\begin{table}[htbp]
\centering
\caption{Controlled Architectural and QoS Ablation Contrasts}
\label{tab:ablation}
\begin{tabular}{@{}llrr@{}}
\toprule
Comparison & Varies & $\Delta\rho$ (mean) & $\Delta\text{F1}$ (mean) \\
\midrule
Topo-QoS $-$ Topo-BL & QoS weighting on structural metrics & \textbf{+0.379} & \textbf{+0.134} \\
HGL $-$ GL & Homogeneous $\to$ heterogeneous architecture & \textbf{+0.225} & \textbf{+0.092} \\
GL-QoS $-$ GL & Scalar QoS edge weight in homogeneous GAT & \textbf{+0.033} & \textbf{+0.010} \\
HGL-QoS $-$ HGL & Explicit QoS vector in heterogeneous GAT & \textbf{$-$0.035} & \textbf{$-$0.004} \\
\bottomrule
\end{tabular}
\end{table}

The ablation analysis reveals that incorporating QoS information into structural centrality metrics
helps substantially (Topo-QoS beats Topo-BL by the largest margin in the table), but the explicit
addition of QoS edge attributes to the heterogeneous GNN does not, on average, improve it
in-distribution --- a small negative effect, consistent with the conference submission's original
finding on this specific point. The architectural contrast remains the larger in-distribution
result: preserving typed node and relation semantics yields an F1 improvement of 0.092 even in the
absence of explicit QoS attributes, suggesting the pub-sub topology and typed relation structure
already capture much of the QoS-relevant routing within a trained scenario.

Unlike the conference submission, we do not find that QoS encoding compensates
out-of-distribution. \S~5.4 shows both HGL and HGL-QoS falling to negative mean rank correlation on
unseen scenarios --- HGL-QoS's LOSO $\rho=-0.058$ is less negative than HGL's $-0.114$, a small
relative improvement in the same direction the conference version reported, but both remain well
below the homogeneous baselines' near-zero LOSO $\rho$ (GL: $+0.021$, GL-QoS: $+0.002$). Explicit
QoS encoding is therefore not the generalization mechanism we previously reported it to be: it
neither rescues out-of-distribution performance in absolute terms, nor does the
in-distribution/out-of-distribution trade-off resolve in HGL-QoS's favor the way a purely relative
reading of ``less negative than HGL'' might suggest. \S~5.4 examines what actually happens under
LOSO, and \S~6.1 discusses why an undisclosed ensemble step in the pipeline that generated the
original conference-submission numbers is the most likely explanation for the discrepancy.

\subsection{Leave-One-Scenario-Out Generalization}

To assess out-of-distribution performance, we employ Leave-One-Scenario-Out (LOSO) cross-validation.
In each iteration, the model is trained on six scenarios and evaluated on the remaining, held-out
scenario. This protocol closely approximates real-world pre-deployment conditions, where a
predictive model must be applied to a new pub-sub architecture with unobserved cascade dynamics.

\begin{table}[htbp]
\centering
\caption{Leave-One-Scenario-Out (LOSO) Cross-Validation Results}
\label{tab:loso}
\begin{tabular}{@{}lrrrr@{}}
\toprule
Variant & Mean $\rho$ & Std $\rho$ & F1@K & $\Delta\rho$ vs GL \\
\midrule
\texttt{GL} & \textbf{0.0208} & 0.1418 & \textbf{0.2086} & --- \\
\texttt{GL-QoS} & 0.0024 & 0.0946 & 0.2008 & $-0.0184$ \\
\texttt{HGL} & $-0.1141$ & 0.1362 & 0.1814 & $-0.1349$ \\
\texttt{HGL-QoS} & $-0.0578$ & 0.1124 & 0.1848 & $-0.0786$ \\
\bottomrule
\end{tabular}
\end{table}

This is the paper's central negative result, and we report it as found. Neither heterogeneous
variant generalizes to a genuinely unseen scenario topology as trained here: HGL's mean
out-of-distribution correlation is $\rho=-0.1141$ and HGL-QoS's is $\rho=-0.0578$ --- both negative,
both below the homogeneous GL ($\rho=0.0208$) and GL-QoS ($\rho=0.0024$), which themselves barely
exceed zero. F1@K follows the same ordering: GL's 0.2086 is the best of the four variants; HGL and
HGL-QoS trail at 0.1814 and 0.1848. Std $\rho$ is large across all four variants (0.11--0.14),
reflecting the seven-fold LOSO protocol's small per-fold sample and the wide spread visible in the
per-fold breakdown (provided in the replication package): individual folds range from strongly
positive to strongly negative for every variant, and the reported means are the average of that
spread, not a stable per-fold guarantee.

Relative to HGL, HGL-QoS's smaller negative mean ($-0.058$ vs.\ $-0.114$) is the one place explicit
QoS encoding still helps directionally, echoing the in-distribution ablation's finding that QoS
contributes little on its own (\S~5.3) --- but ``less negative'' is not ``positive,'' and neither
heterogeneous variant clears the near-zero homogeneous baselines here. Two things are worth being
precise about. First, this result is a genuine reversal of what an in-distribution-only evaluation
(\S~5.1--5.3) would predict: the same typed message-passing architecture that improves ranking and
identification within a trained scenario's distribution does not, on this evidence, transfer that
advantage to an architecturally different held-out scenario. Second, we found this LOSO story only
by rerunning the full evaluation for this revision; the original conference submission reported
HGL-QoS as the strongest out-of-distribution performer ($\rho=0.401$). \S~6.1 traces that
discrepancy to an undisclosed ensemble-blending step present in the pipeline that generated the
original numbers but absent from the architecture actually described in \S~3.2 --- the version
evaluated here. We regard the result in this table, not the conference submission's, as the one
that reflects the method as documented.

\subsection{Reversed-Projection Ablation}
\label{sec:reversed}

Section~\ref{sec:graphrep} argues that a \texttt{DEPENDS\_ON} arrow drawn subscriber $\to$
publisher looks reversed to a middleware-literate reader even though it is not, because it runs
against the direction of data flow. We now test that argument directly rather than only
definitionally: we invert the direction convention used to build the logical dependency projection
and ask whether a structural predictor built on the inverted graph still tracks $I^*(v)$.

\textbf{Metric choice.} Betweenness centrality --- the metric \texttt{Topo-BL} uses elsewhere in
this paper --- was tried first and rejected for this specific test: summed over all ordered node
pairs, betweenness centrality is provably invariant under globally reversing every edge in a graph,
so it cannot show an ablation effect by construction, whatever the true effect of the convention
is. We instead use an \textbf{ancestor-count} proxy, $\text{centrality}(v) = |\text{ancestors}(v)|
/ (|V|-1)$ on the projection --- i.e., how many components transitively depend on $v$ --- which is
exactly the quantity the paper's convention claims predicts blast radius, and which \emph{is}
direction-sensitive: under the correct dependent $\to$ dependency convention it counts $v$'s
dependents (who is hurt if $v$ fails); under the inverted convention it instead counts $v$'s own
dependencies (what $v$ itself relies on), a materially different and, if the convention matters,
materially worse predictor.

\begin{table}[htbp]
\centering
\caption{Reversed-Projection Ablation --- Spearman $\rho$ Against $I^*(v)$, Application Nodes}
\label{tab:reversed}
\begin{tabular}{@{}lrrr@{}}
\toprule
Scenario & $\rho$, forward (paper convention) & $\rho$, reversed & $\Delta\rho$ \\
\midrule
AV System & 0.913 & $-0.893$ & $-1.805$ \\
IoT Smart City & 0.752 & $-0.546$ & $-1.298$ \\
Financial Trading & 0.877 & $-0.877$ & $-1.753$ \\
Healthcare & 0.898 & $-0.809$ & $-1.707$ \\
Hub-and-Spoke & 0.805 & $-0.805$ & $-1.610$ \\
Microservices & 0.798 & $-0.781$ & $-1.579$ \\
Enterprise & 0.832 & $-0.831$ & $-1.663$ \\
\midrule
\textbf{Mean} & \textbf{0.839} & \textbf{$-0.792$} & \textbf{$-1.631$} \\
\bottomrule
\end{tabular}
\end{table}

Inverting the projection does not merely degrade the structural baseline --- it flips its
correlation with ground-truth cascade impact from strongly positive to strongly negative in every
one of the seven scenarios. This is the direct empirical counterpart to
Section~\ref{sec:graphrep}'s definitional argument: the paper's dependent $\to$ dependency
convention is not an arbitrary choice that reviewers happened to misread --- its reversal is
measurably, substantially wrong as a predictor of which components matter. (Ground truth and the
raw structural graph loading in this experiment reuse the exact code path behind every other number
in this paper --- \texttt{saag/simulation/fault\_injector.py} via \texttt{cli/simulate\_graph.py}'s
graph loader --- so only the projection direction varies.)

\subsection{Hardening-Budget Analysis}
\label{sec:hardening}

The metrics in \S~5.1--5.4 are internal to the prediction task: they measure how well a model's
ranking agrees with the simulator's, not what an operator gains by acting on that ranking. We close
that gap with a hardening-budget analysis. Hardening a component is modeled as giving it a hot
replica or failover path, so that its own failure no longer cascades to its dependents --- under
that model, an operator with a budget for $K$ hardening actions wants to know which $K$ components
to harden to eliminate the largest share of total simulated cascade risk. We report that as
\textbf{risk-mass coverage}: $\sum_{v \in \text{top-}K} I^*(v) \,/\, \sum_{v \in \text{Application}}
I^*(v)$, with $K$ fixed by the same rank-matched-binarization rule used for F1 in \S~4.2 ($K =
|\{v : I^*(v) > 0.5\}|$, falling back to 10\% of application nodes when that set is empty).
Selection is compared across three methods: \textbf{HGL} uses the Leave-One-Scenario-Out predicted
score for each scenario (\S~5.4) --- i.e., a model that has never seen this scenario during
training, the realistic condition for a genuinely new pre-deployment system; \textbf{Betweenness}
uses the \S~5.5 ancestor-count centrality computed with full in-scenario visibility;
\textbf{Random} draws $K$ application nodes uniformly (seed 42).

\begin{table}[htbp]
\centering
\caption{Hardening-Budget Risk-Mass Coverage by Top-$K$ Selection Method}
\label{tab:hardening}
\begin{tabular}{@{}lrrrr@{}}
\toprule
Scenario & $K$ & HGL (LOSO) & Betweenness & Random \\
\midrule
AV System & 9 & 1.4\% & 1.8\% & 0.9\% \\
IoT Smart City & 20 & 6.4\% & 21.0\% & 18.4\% \\
Financial Trading & 8 & 0.7\% & 48.8\% & 0.8\% \\
Healthcare & 10 & 39.1\% & 60.6\% & 20.1\% \\
Hub-and-Spoke & 12 & 24.6\% & 16.7\% & 16.8\% \\
Microservices & 34 & 38.3\% & 53.0\% & 21.2\% \\
Enterprise & 9 & 0.3\% & 30.6\% & 0.3\% \\
\midrule
\textbf{Mean} & --- & \textbf{15.8\%} & \textbf{33.2\%} & \textbf{11.2\%} \\
\bottomrule
\end{tabular}
\end{table}

This table now consistently reflects \S~5.4's LOSO finding rather than contradicting it. HGL's
LOSO-selected top-$K$ covers a mean of only 15.8\% of total simulated cascade risk --- barely above
Random's 11.2\%, and less than half of Betweenness's 33.2\%. This is the same negative
out-of-distribution result as Table~\ref{tab:loso}, restated in an operator-facing metric: a
component-hardening budget allocated by a LOSO-evaluated HGL model, as trained here, would capture
materially less cascade risk than the same budget allocated by a cheap, deterministic structural
heuristic with full visibility into the target scenario. HGL's coverage still exceeds Random in
five of seven scenarios (all but IoT Smart City and Financial Trading, where it falls to single
digits), so LOSO-HGL predictions are not pure noise --- but they are a materially worse guide to
hardening decisions than either the structural baseline or, per \S~5.1--5.3, an
in-distribution-trained HGL model would be. A like-for-like comparison --- in-distribution HGL
predictions (matching the regime behind this paper's in-distribution F1 = 0.688) against
in-distribution betweenness --- is left for future work; it requires GNN inference from the trained
per-scenario checkpoints rather than the LOSO predictions used here.

\emph{Methodological note.} An initial design for this experiment measured hardening's effect by
physically removing the top-$K$ nodes from the topology and re-running the fault-injection engine
on the remaining graph. That design was discarded after its own output invalidated it: it made
every selection method's ``hardening'' look harmful (mean impact on remaining nodes
\emph{increased}), because deleting a publisher node destroys it as a message source for its
subscribers permanently --- the opposite of hardening, which keeps the component functioning while
removing its failure risk. We flag this here because the failure was informative: it is a concrete
illustration of why simulator semantics (\S~3) have to be gotten right before a derived experiment
can be trusted, which is the same lesson Section~\ref{sec:graphrep} and this section's own
risk-mass framing are built around.

% ===========================================================================
\section{Threats to Validity}

We evaluate threats to validity along three dimensions: construct, internal, and external
validity. The most notable limitation of this study is that our ground-truth labels are generated
via a controlled simulation rather than derived from observed failures in real-world
publish-subscribe systems. As such, the strongest claims we can make are comparative: our results
speak to which modeling choices perform better under identical experimental conditions, rather
than providing guarantees of absolute predictive accuracy in operational deployments.

\subsection{Predictor Reproducibility: An Undisclosed Ensemble Step in the Conference-Submission Numbers}
\label{sec:ensemble}

Every number reported in \S~5 of this revision comes from a full rerun of the evaluation pipeline
performed for this submission, at the 300-epoch configuration \S~3.3 describes. That rerun did not
reproduce the conference submission's numbers: in-distribution ranking correlation and F1 dropped
for every learned variant relative to the conference submission (\S~5.1--5.3), and the
Leave-One-Scenario-Out result reversed sign for HGL and HGL-QoS (\S~5.4). Structural baselines
(Topo-BL, Topo-QoS), which involve no learned model, reproduced exactly --- ruling out a change in
the ground-truth simulation or the underlying scenario data as the cause.

We traced the discrepancy rather than reporting it as unexplained variance. A single fixed
seed/scenario/variant cell (Hub-and-Spoke, HGL, seed 42) that scored $\rho=0.9605$ in the cached
results behind the conference submission scored $\rho=-0.2188$ under this revision's pipeline, run
three times, bit-identically. This ruled out both non-determinism and the epoch count (300 vs.\ the
originally-run 200; both epoch settings reproduced $\rho=-0.2188$ for that cell) as explanations,
leaving a genuine divergence between the code that generated the cached numbers and the code
evaluated here. Inspecting the commit history identified the cause: the pipeline that generated the
conference submission's numbers computed a blended score
$Q_{\text{ens}}(v) = \alpha \cdot Q_{\text{GNN}}(v) + (1-\alpha) \cdot Q_{\text{RMAV}}(v)$, combining
the GNN's output with a separate, non-learned multi-dimensional quality-attribution score
($Q_{\text{RMAV}}$, part of a related but distinct framework not otherwise described in this paper)
before evaluating against $I^*(v)$. This ensemble step was removed from the codebase during later,
unrelated refactoring; the current pipeline evaluates the pure GNN score $Q_{\text{GNN}}(v)$ that
\S~3.2 actually describes (``HGL is realized as a heterogeneous graph attention network...'' --- no
ensemble, no RMAV blending appears anywhere in that description).

We regard this revision's numbers, not the conference submission's, as the ones that validly
support the paper's methodology section, precisely because they are what a reader implementing
\S~3.2 as written would obtain. The blended score's better performance --- particularly its much
stronger apparent out-of-distribution generalization --- is itself informative: an RMAV-style
structural composite is not learned per-scenario, so blending it in would mechanically stabilize
LOSO predictions regardless of what the GNN component contributes, which is consistent with the
sharp LOSO-specific reversal we observe once the blend is removed (\S~5.4). We deliberately do not
reintroduce that ensemble step here, disclosed or otherwise: doing so would substitute a second,
non-learned predictor for part of what this paper claims the GNN itself achieves, which is the
exact methodological ambiguity this section exists to flag and correct. HGL in this paper is, and
remains, the pure heterogeneous graph attention network \S~3.2 describes --- nothing else
contributes to $Q^*(v)$. \S~7 discusses improving its out-of-distribution behavior on those terms.

\subsection{Simulator-Derived Labels}

The target impact score $I^*(v)$ is calculated using a discrete-event cascade simulator that
operates on the same graph-based system model used by the learning algorithms. Consequently, both
the predictions $Q^*(v)$ and the targets $I^*(v)$ are linked to a shared scenario specification,
but are produced by different processes: HGL uses graph features and relation-specific message
passing, while the simulator explicitly propagates faults to compute cascade impact. This approach
introduces the risk of circular validation. High concordance with $I^*(v)$ indicates that a model
has learned the simulator's notion of cascade semantics, but does not necessarily imply identical
behavior in real or production environments.

To address this potential bias, we have implemented several safeguards. First, all model variants
are evaluated against the same simulator-derived targets, ensuring that comparative performance
among HGL, homogeneous GNNs, and structural baselines is internally consistent. Second,
ground-truth labels are generated using a softened dynamic simulation rather than static structural
proxies, thereby minimizing direct overlap between test targets and input features. Third, all
learned models are tested in a feature-decoupled setting where pre-computed centrality metrics are
withheld as input. While these measures enhance the validity of our architectural comparisons, we
underscore the need for additional external validation (ideally using measured runtime failures) to
support strong claims regarding real-world deployment accuracy.

\subsection{Per-Type Label Informativeness}

Because the simulator models cascade impact as degradation in pub-sub message flows, Library nodes
consistently receive near-constant, close-to-zero impact scores: failures in shared libraries or
dependencies are not simulated as cascade-triggering events. As a result, there is almost no
variance among Library node labels, and their within-type Spearman correlation is essentially
undefined ($\approx 0.000$).

This design choice has two important implications for interpreting our results. First, the core
predictive target in this study is \emph{Application-level} criticality, and headline correlation
measures should be understood as reflecting Application nodes. Aggregated statistics across all
node types may obscure this fact, as they combine the informative Application-node distribution
with the degenerate Library-node one---an instance of Simpson's paradox
\cite{Simpson1951,Bickel1975}, in which a meaningful per-type signal is masked in the overall
average. We therefore treat per-type, stratified reporting as a methodological contribution in its
own right, and base our headline claims on Application-node performance rather than on the pooled
aggregate. Second, these limitations point to a clear extension: enhancing the simulator so that
library failures propagate to all dependent applications would make Library nodes authentic
prediction targets. This would be a setting where heterogeneous, type-aware modeling is likely to
yield its greatest performance benefits over homogeneous baselines. Rather than omitting Library
nodes---an action that would forfeit the very heterogeneity our architecture is designed to
leverage---we retain them in our models, but restrict any claims about Library-node predictions to
the architectural level, leaving behavioral validation of library-failure propagation for future
work.

\subsection{Training and Ablation Design}

Our GNN training configuration is standardized across all model variants: each model employs 4
attention heads per relation, a hidden dimension of 64, 300 training epochs, the AdamW optimizer
with a learning rate of $10^{-3}$, a dropout rate of 0.2, and weight decay of $10^{-4}$. This fixed
setup ensures a fair basis for comparison, though it may not be optimal for every variant. For
example, the HGL-QoS variant---with its additional QoS edge-feature dimensions---could benefit from
different choices of model capacity, regularization strategies, or learning-rate schedules.

To partially address this concern, we conducted targeted sensitivity analyses, varying the learning
rates ($5\times10^{-4}$, $10^{-3}$, $2\times10^{-3}$) and hidden dimensions ($64, 128$) in
scenarios where HGL-QoS initially underperformed HGL. Across all configurations, the qualitative
results remained robust: explicit QoS edge-feature injection did not consistently outperform the
QoS-masked HGL in in-distribution evaluations. While this does not preclude future improvements
through broader hyperparameter or architecture searches, it suggests that our main findings are not
the product of a single suboptimal training setting.

A related sensitivity surface that remains unquantified concerns the simulation-softening
constants. The QoS-factor multipliers applied during cascade softening are fixed point values:
$\times1.2$ for RELIABLE topics, $\times1.15$ for HIGH/CRITICAL/URGENT transport priority, and
$\times1.05$ for MEDIUM priority. Like the propagation threshold $\theta$ and the AHP shrinkage
factor $\lambda$, these constants affect the continuous target $I^*(v)$ and thus the ranking labels
against which all models are evaluated. Their influence on rank stability and the relative ordering
of model variants has not been separately quantified; a systematic sensitivity analysis across
these simulation hyperparameters is left as future work.

A further consideration regarding internal validity is label sparsity. In decoupled pub-sub
topologies, raw fault-injection simulations often produce many zero-impact labels, which can bias
models toward degenerate, constant predictions. To alleviate this, we employ a Simulation Softening
procedure that produces continuous cascade-impact targets based on rate-weighted failed-publisher
fractions and topic-level QoS factors. This method enables more meaningful learning and ranking
evaluations, though it also means that the target reflects a softened cascade process rather than
strict binary failure propagation.

\subsection{Scenario Coverage and Scale}

To assess external validity, our evaluation encompasses seven diverse scenarios, including
autonomous vehicles, financial trading, healthcare integration, smart-city telemetry,
hub-and-spoke enterprise integration, cloud-native microservices, and large-scale enterprise
pub-sub systems. This variety captures a broad spectrum of topology densities, broker fan-outs,
QoS heterogeneity, and application counts. Nevertheless, it is important to note that all scenarios
are synthetic, generated from parameterized configurations rather than from actual production
deployments. Real-world systems may exhibit different QoS distributions, correlated failure
patterns, operational interventions, workload variability, and deployment constraints---factors
that are not fully reflected in our synthetic evaluation.

There are also limitations regarding scale. Most of the scenarios analyzed in this study contain on
the order of tens of Application and Library nodes; the largest scenario, the \texttt{xlarge}
Enterprise deployment, comprises on the order of 300 applications. As a result, our finding that
explicit QoS features do not enhance in-distribution HGL performance may not generalize to
substantially larger systems, those with many hundreds to thousands of application-level nodes,
where greater data availability could enable more effective learning of QoS-dependent aggregation.
Thus, we confine this claim to the evaluated scales and scenario types. Evaluating substantially
larger topologies, those comprising several thousand application-level nodes, remains an open
avenue for future research and is necessary before extending these conclusions to true large-scale
deployments.

Finally, while HGL carries out message passing across the entire multi-tier graph, our
ground-truth labels reflect cascade impact only at the application level: the simulator quantifies
message-flow degradation among applications but does not attribute independent impact to Topic,
Broker, or compute Node tiers. Consequently, although these tiers are part of the learned
representation, our evaluation focuses on predictive performance for Application (and Library)
nodes, and we make no validated claims about infrastructure-level criticality. Enhancing the
simulator to assign cascade impact to infrastructure tiers---so that broker or node failures become
explicit prediction targets---is a clear direction for further strengthening these results.

\paragraph{Remaining Experimental Work.} The reversed-projection ablation and hardening-budget
experiment called for earlier in this revision are now reported in \S~5.5 and \S~5.6, and both
were rerun for this revision alongside the rest of \S~5 (\S~6.1). One follow-up they surface
remains open: \S~5.6's hardening-budget comparison holds HGL to its harder Leave-One-Scenario-Out
condition while giving the structural baseline full in-distribution visibility, which is not a
like-for-like comparison. Repeating it with in-distribution HGL predictions --- matching the regime
behind this paper's in-distribution F1 = 0.688 (\S~5.2) --- requires running inference from the
trained per-scenario checkpoints (\texttt{output/gnn\_checkpoints/}) rather than reusing the
already-available LOSO predictions, and is scoped as a follow-up before submission.

% ===========================================================================
\section{Conclusion}

This paper presents HGL, a heterogeneous graph learning framework for predicting the impact of
component-level failure cascades using pre-deployment graph-based models of publish-subscribe
middleware. By applying relation-aware message passing over the native heterogeneous graph---which
encompasses Applications, Libraries, Topics, Brokers, and deployment Nodes---HGL preserves the
typed semantics that are lost when systems are reduced to single-type logical dependency
projections, as in homogeneous baselines.

Across seven representative pub-sub scenarios and within a rigorously controlled $2\times3$
experimental design, our findings establish typed heterogeneous graph learning as a real but
conditional predictive advantage. Within a trained scenario's distribution, HGL achieves the
highest mean ranking performance among the learned variants (Spearman $\rho=0.411$) and improves
critical-component identification over homogeneous graph learning (mean F1 = 0.688,
$\Delta\text{F1}=+0.092$ over GL), though a QoS-weighted structural baseline (Topo-QoS) remains
competitive or superior in several scenarios (\S~5.1--5.2). Ablation studies show these
in-distribution gains come principally from modeling heterogeneous node and relation semantics
rather than from explicit QoS edge features (\S~5.3). Out-of-distribution, the picture reverses:
Leave-One-Scenario-Out validation shows both HGL and HGL-QoS falling to negative mean rank
correlation on unseen scenarios, underperforming the near-zero homogeneous baselines (\S~5.4). We
report this reversal as found --- it is the paper's central negative result, and it directly
contradicts what an in-distribution-only evaluation would suggest.

That reversal, together with a substantial in-distribution numerical gap from the conference
submission this paper extends, has a traceable cause rather than being unexplained noise: the
evaluation pipeline behind the conference submission's numbers computed a blend of the GNN's score
with a separate, non-learned quality-attribution score, an ensemble step never described in this
paper's methodology and since removed from the codebase (\S~6.1). The numbers in this revision
reflect the pure heterogeneous-GNN architecture \S~3.2 actually describes. We surface this
transparently, in the same spirit as this paper's other honestly-reported negative findings --- the
Simpson's-paradox masking of Library-node signal \cite{Simpson1951,Bickel1975} and the
hardening-budget experiment's own discarded first design (\S~5.6) --- because a paper whose central
claim rests on a rerun that contradicts the original submission is exactly the situation
methodological transparency exists for.

By natively modeling the complete five-type architecture (Applications, Libraries, Topics,
Brokers, and deployment Nodes), HGL uses typed relations to retain semantics that homogeneous
Graph Neural Networks (GNNs) overlook --- within a trained scenario. Importantly, while HGL
incorporates this multi-tier infrastructure to provide a rich contextual embedding layer, our
current predictions and empirical validations focus specifically on cascade impacts at the
application level.

We acknowledge several important limitations of this study, foremost among them the ones detailed
above. The reported targets are generated via simulation, and the evaluated scenarios are
synthetic; consequently, absolute metric values should be interpreted as simulator-based estimates
rather than guarantees of operational accuracy. Our most defensible contribution is therefore
narrower than the conference submission claimed: within the disclosed evaluation framework,
preserving heterogeneous middleware semantics yields better pre-deployment identification of
critical components than homogeneous graph learning, but only within a trained scenario's
distribution --- generalization to unseen topologies, as evaluated here, is not yet achieved by
either the heterogeneous or the homogeneous learned variants, and the deterministic structural
baseline remains the more reliable out-of-distribution choice on this evidence (\S~5.6).

In the future, we will extend this work in four directions. First, \S~6.1's finding motivates
improving HGL's out-of-distribution robustness directly, as a single learned architecture, rather
than by reintroducing a second, non-learned predictor via ensembling --- a direction we
deliberately do not pursue, since it would reopen the exact methodological-transparency problem
\S~6.1 diagnoses. Promising directions on those terms include training over a larger and more
topologically diverse scenario corpus, explicit domain-generalization regularization, or
meta-learning objectives that directly optimize for cross-scenario transfer rather than
in-distribution fit. Second, although HGL currently performs message passing across the full
multi-tier graph, ground-truth impact scores are limited to the application level; expanding the
simulator to attribute cascade impact to Topics, Brokers, and compute Nodes would enable direct
validation of HGL's predictions for these infrastructure layers. Third, we aim to validate HGL
against real-world failure data from operational deployments such as Kubernetes- or ROS2-based
publish-subscribe systems, essential to move beyond simulator-based evidence. Fourth, a systematic
comparison against other heterogeneous architectures --- such as RGCN~\cite{Schlichtkrull2018},
HAN~\cite{Wang2019}, HGT~\cite{Hu2020}, and MAGNN~\cite{Fu2020} --- would help establish whether the
out-of-distribution weakness observed here is specific to this HeteroGAT instantiation or a broader
property of typed message passing under scenario shift.

\section*{Acknowledgment}

During the preparation of this work, the author(s) used Gemini Flash 3.5 and Grammarly in order to
improve language and readability. After using this tool, the author(s) reviewed and edited the
content as needed and take(s) full responsibility for the content of the publication.

\begin{thebibliography}{99}

\bibitem{Eugster2003} P. T. Eugster, P. A. Felber, R. Guerraoui, A.-M. Kermarrec, ``The many faces
of publish/subscribe,'' \emph{ACM Computing Surveys}, vol. 35, no. 2, pp. 114--131, 2003.

\bibitem{Carzaniga2001} A. Carzaniga, D. S. Rosenblum, A. L. Wolf, ``Design and evaluation of a
wide-area event notification service,'' \emph{ACM Transactions on Computer Systems}, vol. 19,
no. 3, pp. 332--383, 2001.

\bibitem{OMG2015DDS} Object Management Group, \emph{Data Distribution Service (DDS), Version 1.4},
OMG formal specification, 2015.

\bibitem{OASIS2019MQTT} A. Banks, E. Briggs, K. Borgendale, R. Gupta, \emph{MQTT Version 5.0},
OASIS Standard, 2019.

\bibitem{Freeman1977} L. C. Freeman, ``A set of measures of centrality based on betweenness,''
\emph{Sociometry}, vol. 40, no. 1, pp. 35--41, 1977.

\bibitem{Brandes2001} U. Brandes, ``A faster algorithm for betweenness centrality,'' \emph{Journal
of Mathematical Sociology}, vol. 25, no. 2, pp. 163--177, 2001.

\bibitem{BrinPage1998} S. Brin, L. Page, ``The anatomy of a large-scale hypertextual web search
engine,'' \emph{Computer Networks and ISDN Systems}, vol. 30, no. 1--7, pp. 107--117, 1998.

\bibitem{Albert2000} R. Albert, H. Jeong, A.-L. Barab\'{a}si, ``Error and attack tolerance of
complex networks,'' \emph{Nature}, vol. 406, pp. 378--382, 2000.

\bibitem{Buldyrev2010} S. V. Buldyrev, R. Parshani, G. Paul, H. E. Stanley, S. Havlin,
``Catastrophic cascade of failures in interdependent networks,'' \emph{Nature}, vol. 464,
pp. 1025--1028, 2010.

\bibitem{Kipf2017} T. N. Kipf, M. Welling, ``Semi-supervised classification with graph
convolutional networks,'' in \emph{International Conference on Learning Representations (ICLR)},
2017.

\bibitem{Hamilton2017} W. L. Hamilton, R. Ying, J. Leskovec, ``Inductive representation learning on
large graphs,'' in \emph{Advances in Neural Information Processing Systems (NeurIPS)}, pp.
1024--1034, 2017.

\bibitem{Velickovic2018} P. Veli\v{c}kovi\'{c}, G. Cucurull, A. Casanova, A. Romero, P. Li\`{o},
Y. Bengio, ``Graph Attention Networks,'' in \emph{International Conference on Learning
Representations (ICLR)}, 2018.

\bibitem{Fan2020} J. Fan, L. Wang, J. Liu, ``FINDER: Free-form Information Diffusion and Networked
Evaluation in Graphs,'' \emph{Nature Machine Intelligence}, vol. 2, no. 6, pp. 338--346, 2020.

\bibitem{Munikoti2022} S. Munikoti, D. Agarwal, L. de Melo, ``DrBC: Betweenness Centrality
Estimation using Graph Convolutional Networks,'' \emph{Neurocomputing}, vol. 491, pp. 215--226,
2022.

\bibitem{Munikoti2024} S. Munikoti et al., ``PowerGraph: Using Graph Neural Networks for Critical
Node Identification in Power Grids,'' in \emph{Advances in Neural Information Processing Systems
(NeurIPS)}, 2024.

\bibitem{Schlichtkrull2018} M. Schlichtkrull, T. N. Kipf, P. Bloem, R. van den Berg, I. Titov,
M. Welling, ``Modeling Relational Data with Graph Convolutional Networks,'' in \emph{Extended
Semantic Web Conference (ESWC)}, pp. 593--607, Springer, 2018.

\bibitem{Wang2019} X. Wang, H. Ji, C. Shi, B. Wang, Y. Ye, P. Cui, P. S. Yu, ``Heterogeneous Graph
Attention Network,'' in \emph{Proceedings of the Web Conference (WWW)}, pp. 2022--2032, 2019.

\bibitem{Hu2020} Z. Hu, Y. Dong, K. Wang, Y. Sun, ``Heterogeneous Graph Transformer,'' in
\emph{Proceedings of the Web Conference (WWW)}, pp. 2704--2710, 2020.

\bibitem{Fu2020} X. Fu, J. Zhang, Z. Meng, C. Shi, ``MAGNN: Metapath-aggregated Graph Neural
Network for Heterogeneous Graphs,'' in \emph{Proceedings of the Web Conference (WWW)}, pp.
2331--2341, 2020.

\bibitem{Li2018} Q. Li, Z. Han, X.-M. Wu, ``Deeper insights into graph convolutional networks for
semi-supervised learning,'' in \emph{AAAI Conference on Artificial Intelligence}, 2018.

\bibitem{Simpson1951} E. H. Simpson, ``The interpretation of interaction in contingency tables,''
\emph{Journal of the Royal Statistical Society: Series B (Methodological)}, vol. 13, no. 2, pp.
238--241, 1951.

\bibitem{Bickel1975} P. J. Bickel, E. A. Hammel, J. W. O'Connell, ``Sex bias in graduate
admissions: Data from Berkeley,'' \emph{Science}, vol. 187, no. 4175, pp. 398--404, 1975.

\end{thebibliography}

\end{document}
